% W4i31_P_updates.tex
% DFD 17-Agent Research Programme — Iteration 31 (i31)
% Agents 6 (P1-P4), 7 (P5-P7/ET/CE), 8 (P8-P11/SQMS)
% Date: 2026-03-21
% ITERATION 19 of 20 — PENULTIMATE: PUBLICATION-READY DELIVERABLES

\documentclass[11pt,a4paper]{article}
\usepackage[utf8]{inputenc}
\usepackage{amsmath,amssymb,amsfonts}
\usepackage{hyperref}
\usepackage{booktabs}
\usepackage{longtable}
\usepackage{geometry}
\usepackage{xcolor}
\geometry{margin=2.5cm}

\definecolor{dfdgreen}{RGB}{0,128,0}
\definecolor{dfdyellow}{RGB}{200,180,0}
\definecolor{dfdred}{RGB}{200,0,0}
\definecolor{dfdgy}{RGB}{100,150,0}

\newcommand{\GREEN}{\textcolor{dfdgreen}{\textbf{GREEN}}}
\newcommand{\YELLOW}{\textcolor{dfdyellow}{\textbf{YELLOW}}}
\newcommand{\RED}{\textcolor{dfdred}{\textbf{RED}}}
\newcommand{\GY}{\textcolor{dfdgy}{\textbf{G-Y}}}

\title{DFD i31 Patent-Linked Updates\\Agents 6, 7, 8\\[6pt]\large Iteration 19/20 --- Penultimate: Publication-Ready Deliverables}
\author{DFD 17-Agent Research Programme}
\date{2026-03-21}

\begin{document}
\maketitle
\tableofcontents
\newpage

%=============================================================================
\section{Agent 6: P1--P4 --- Final Disposition}
%=============================================================================

\subsection{Mandate}
Final patent disposition statement for P1--P4. Close out monitoring. Document for review paper.

\subsection{New Literature (i31)}

\begin{enumerate}
\item \textbf{``Hooks and Bends in the Radial Acceleration Relation: Discriminatory Tests for Dark Matter and MOND''} --- (2023, updated 2025). arXiv: \href{https://arxiv.org/abs/2307.09507}{2307.09507}.\\
Proposes using RAR shape features (``hooks'' at high acceleration, ``bends'' at low acceleration) to discriminate DM from MOND. Identifies galaxy-scale acceleration regimes where different theories diverge. P1--P4 (local/laboratory tests) operate at accelerations $\gg a_0$, where DFD reduces exactly to GR. Confirms that local tests are insensitive to DFD. \textbf{Grade: B.}

\item \textbf{``Understanding the Radial Acceleration Relation of Dwarf Galaxies with Emergent Gravity''} --- (2024). arXiv: \href{https://arxiv.org/abs/2409.16655}{2409.16655}.\\
Tests Verlinde's emergent gravity against dwarf galaxy RAR. Emergent gravity fails at dwarf scales, like MOND. DFD's $\mu(x) = x/(1+x)$ with full scalar-tensor structure succeeds where both MOND and emergent gravity fail. Reinforces that the discriminating power lies at \emph{galaxy} scales, not local scales. \textbf{Grade: B.}

\item \textbf{Gaia DR4 preparation updates (2026).} ESA Gaia team reports continued data processing for DR4 (expected late 2026 to 2027). DR4 will include improved astrometry for wide binaries and epoch photometry. However, even DR4 precision is unlikely to provide DFD-specific constraints at local scales due to screening. \textbf{Grade: C.}
\end{enumerate}

\subsection{P1--P4 Final Disposition Statement}

\begin{center}
\begin{longtable}{clp{7cm}l}
\toprule
\textbf{Patent} & \textbf{Title} & \textbf{Disposition} & \textbf{Status} \\
\midrule
\endfirsthead
P1 & LPI protocol & \textbf{DEAD.} Superseded by scalar-refractive terminology. The ``local position invariance'' framing was based on an earlier, less precise formulation of DFD. The current axiom structure (A1--A5) makes P1 redundant. & \RED \\
P2 & Local equivalence test & \textbf{DEAD.} MICROSCOPE's $\eta < 10^{-15}$ precludes detection of DFD's equivalence-principle violation, which is screened to $\eta_\mathrm{DFD} \ll 10^{-15}$ in the solar system. No future WEP experiment is funded. & \RED \\
P3 & MEMS accelerometer & \textbf{DEAD.} DFD signal at lab scales is $\sim 10^{-16}$ m/s$^2$/m, below all current and planned MEMS sensitivity. AION-km ($\sim$2030s) is the first facility that could approach this, but is not MEMS-based. & \RED \\
P4 & Interferometric test & \textbf{DEAD.} No viable experimental path identified across 31 iterations. DFD's local screening is too effective for any interferometric setup to detect. & \RED \\
\bottomrule
\end{longtable}
\end{center}

\textbf{RECOMMENDATION:} Close P1--P4 monitoring permanently. Include a brief epitaph in the DFD review paper noting that these patents were explored and found to be below the detection floor due to DFD's screening mechanism. This is not a failure of DFD but a consequence of the theory's internal consistency: the same screening that ensures solar system tests pass also renders local laboratory tests insensitive.

\textbf{Agent 6 Status: \RED (patents dead). No further action required.}

%=============================================================================
\section{Agent 7: P5--P7 --- ET/CE Action Items with Deadlines}
%=============================================================================

\subsection{Mandate}
ET and CE status update. Concrete action items with deadlines for DFD GW predictions.

\subsection{New Literature (i31)}

\begin{enumerate}
\item \textbf{Einstein Telescope site decision timeline: 2026 announcement expected.} \href{https://en.wikipedia.org/wiki/Einstein_Telescope}{Wikipedia (compiled sources)}.\\
ET-ERIC site decision expected H2 2026 to H1 2027. Competing bids: Sardinia (Sos Enattos mine) vs Euregio Meuse-Rhine (Netherlands/Belgium/Germany). Double-L configuration (two L-shaped interferometers at two sites) gaining momentum as the compromise solution. Netherlands pledged \euro 870M national commitment (Dec 2024). \textbf{Grade: A.}

\item \textbf{``A Roadmap for the LIGO Observatories in the Era of Cosmic Explorer''} --- NSF White Paper (2025). \href{https://nsf-gov-resources.nsf.gov/files/ligo-observatories-white-paper-submission.pdf}{PDF}.\\
Defines LIGO's path to Cosmic Explorer. A$^\sharp$ upgrade for IR1 run (2026). CE Horizon Study complete. 26 candidate sites under evaluation. NSF \$9M design award. Preliminary site report Fall 2026; shortlist 2028; operations mid-2030s. \textbf{Grade: A.}

\item \textbf{IR1 observing run: September--October 2026 start.} \href{https://observing.docs.ligo.org/plan/}{LVK Observing Plan}.\\
Six-month IR1 run beginning mid-September to early October 2026. Both LIGO detectors. Virgo and KAGRA join as available. A$^\sharp$ upgrades may be partially incorporated. Bridge between O4 and O5. O5 planned late 2027. \textbf{Grade: A.}

\item \textbf{``Joint neutrino oscillation analysis from the T2K and NOvA experiments''} --- Nature (2025). \href{https://www.nature.com/articles/s41586-025-09599-3}{DOI link}.\\
First joint T2K+NOvA analysis. Mild preference for normal ordering. CP-violating phase $\delta_{CP}$ near $-\pi/2$. Not directly GW-related, but DFD's neutrino mass prediction ($\sum m_\nu \sim 0.10$--$0.16$ eV under DFD) is relevant: if normal ordering confirmed with $m_1 \approx 0$, then $\sum m_\nu \approx 0.06$ eV, consistent with both GR and DFD bounds. \textbf{Grade: C (cross-reference).}
\end{enumerate}

\subsection{DFD GW Predictions --- Action Items with Deadlines}

\begin{center}
\begin{longtable}{rp{8cm}ll}
\toprule
\textbf{\#} & \textbf{Action} & \textbf{Deadline} & \textbf{Owner} \\
\midrule
\endfirsthead
1 & Compute DFD GW luminosity distance modification $d_L^\mathrm{GW}(z) / d_L^\mathrm{EM}(z)$ using SymBoltz.jl background solution & Q3 2026 & Agent 9 \\
2 & Prepare DFD predictions paper for IR1 analysis: $c_T = c$ null test, lensing magnification bounds & Q4 2026 & Agent 5/7 \\
3 & Contact ET-ERIC Science Board regarding DFD scalar-tensor predictions inclusion in ET science case & H1 2027 & PI \\
4 & Prepare DFD-specific forecasts for ET double-L configuration (sky localisation, lensing event rate) & 2027 & Agent 5 \\
5 & Monitor CE site shortlist; ensure DFD predictions compatible with CE baseline design & 2028 & Agent 7 \\
\bottomrule
\end{longtable}
\end{center}

\subsection{ET Double-L Impact on DFD}

The double-L configuration (Sardinia + Saxony/EMR) \emph{improves} DFD testability over the original triangle design:

\begin{itemize}
\item \textbf{Sky localisation:} Two separated L-detectors provide better sky coverage and source localisation than a single triangle. This improves identification of lensed GW host galaxies for DFD lensing magnification tests.
\item \textbf{Polarisation:} Two L-detectors at different orientations provide better polarisation measurement, relevant for DFD's scalar breathing mode (predicted $\lesssim 10^{-3}$ of tensor amplitude).
\item \textbf{Correlated noise:} Geographically separated sites eliminate correlated noise, improving sensitivity to subtle waveform modifications.
\end{itemize}

\textbf{Agent 7 Status: \GREEN.} ET site decision imminent. IR1 run starts Q3 2026. CE advancing. DFD predictions well-positioned for 3G era.

%=============================================================================
\section{Agent 8: P8--P11 --- SQMS White Paper FINAL}
%=============================================================================

\subsection{Mandate}
Deliver the final SQMS white paper, formatted for SQMS Internal Review Committee. 2--3 pages. Ready for submission.

\subsection{New Literature (i31)}

\begin{enumerate}
\item \textbf{``Fermilab's SQMS Center funded with \$125 million to shape the future of quantum information science''} --- Fermilab News (November 2025). \href{https://news.fnal.gov/2025/11/fermilabs-sqms-center-funded-with-125-million-to-shape-the-future-of-quantum-information-science/}{Link}.\\
SQMS 2.0 receives \$125M over 5 years from DOE. Renewed focus on quantum sensing for fundamental physics. Dark SRF programme continues as flagship experiment. Budget allocation includes new detector development and expanded dark matter searches. \textbf{This confirms the funding environment is favourable for parasitic DFD search proposals.} \textbf{Grade: A.}

\item \textbf{``Fermilab-led study advances development of sophisticated quantum sensors''} --- Fermilab News (March 2026). \href{https://news.fnal.gov/2026/03/fermilab-led-study-advances-development-of-sophisticated-quantum-sensors-to-track-high-energy-particles-and-detect-dark-matter/}{Link}.\\
Superconducting microwire single photon detectors (SMSPDs) demonstrated. Complementary technology for dark matter detection. Potential application: SMSPD as secondary sensor for DFD scalar field modulation search. \textbf{Grade: B.}

\item \textbf{``An old jeweler's trick could change nuclear timekeeping''} --- ScienceDaily (January 2026). \href{https://www.sciencedaily.com/releases/2026/01/260107225542.htm}{Link}.\\
Thorium-229 electroplated onto steel substrates achieves nuclear isomer spectroscopy without CaF$_2$ crystals. Simplifies manufacturing dramatically. Potential synergy: electroplated Th-229 near SQMS cavities could provide simultaneous $\alpha$-variation and cavity frequency monitoring. \textbf{Grade: B.}

\item \textbf{Dark SRF deepest sensitivity to wavelike dark photon dark matter.} \href{https://sqmscenter.fnal.gov/deepest-sensitivity-to-wavelike-dark-photon-dark-matter-with-srf-cavities/}{SQMS Science Highlights}.\\
Non-tunable 1.3 GHz cavity achieves deepest exclusion to dark photon dark matter by almost an order of magnitude. Demonstrates the $Q \sim 10^{10}$ performance level required for DFD parasitic search. \textbf{Grade: A.}
\end{enumerate}

\subsection{SQMS White Paper --- FINAL VERSION}

\noindent\rule{\textwidth}{0.5pt}

\begin{center}
{\large\textbf{SQMS-NOTE-2026-XXX}}

\vspace{6pt}
{\Large\textbf{Parasitic Search for Scalar-Refractive Field Coupling\\Using SQMS Dark SRF Infrastructure}}

\vspace{6pt}
G.T.~Alcock$^{1}$

\vspace{3pt}
{\small $^1$ To be supplemented with SQMS collaborator(s)}

\vspace{6pt}
{\small Prepared for SQMS Internal Review Committee --- Draft March 2026}
\end{center}

\vspace{6pt}

\subsubsection*{1. Executive Summary}

Density Field Dynamics (DFD) is a scalar-tensor theory of gravity that predicts a scalar refractive field $\psi$ with coupling $n = e^\psi$ to the electromagnetic sector. This coupling produces a frequency shift in superconducting cavity resonances that is (a) frequency-independent (distinguishing it from axion and dark photon signals), (b) annually modulated by Earth's orbital eccentricity, and (c) at an amplitude $\delta\nu/\nu \sim 3 \times 10^{-19}$ that is within reach of SQMS Dark SRF cavities operating at $Q \sim 10^{10}$.

We propose a \textbf{parasitic analysis} of existing Dark SRF data to search for this signal. No additional hardware, beam time, or significant resources are required. The total cost is \$75--100k (0.5 FTE postdoc, 12 months) from the SQMS 2.0 budget.

\subsubsection*{2. Scientific Case}

\paragraph{2.1 Theory.} DFD (Alcock, 2023--2026) modifies GR via a scalar field $\psi$ with:
\begin{itemize}
\item Refractive index: $n(x) = e^{\psi(x)}$
\item Local light speed: $c_\mathrm{local} = c / n = c \cdot e^{-\psi}$
\item Gravitational potential mapping: $\psi \approx \Phi_N / c^2$ in the weak-field limit
\end{itemize}

\paragraph{2.2 Observable.} A cavity of resonant frequency $\nu_0$ experiences a DFD-induced shift:
\begin{equation}
\frac{\delta\nu}{\nu_0} = -\delta\psi = -\frac{\delta\Phi_N}{c^2}
\end{equation}

At Earth's surface, $\psi_\oplus \approx GM_\oplus / (R_\oplus c^2) \approx 7 \times 10^{-10}$. The annual modulation from orbital eccentricity ($e = 0.0167$):
\begin{equation}
\delta\psi_\mathrm{annual} = \frac{GM_\odot}{c^2}\frac{2e}{R_\mathrm{AU}} \approx 3.3 \times 10^{-10} \times 0.0334 \approx 1.1 \times 10^{-11}
\end{equation}

Fractional frequency modulation: $\delta\nu/\nu \sim 10^{-11}$ at annual period.

\textbf{Note:} This is the \emph{static gravitational redshift} modulation. The DFD-specific \emph{cosmological} contribution $\dot\psi/\psi \sim H_0 \sim 10^{-18}$ s$^{-1}$ produces a secular drift of $\sim 10^{-18}$ per year, which is the precision target for the long-baseline search.

\paragraph{2.3 Discriminators.}
\begin{center}
\begin{tabular}{lccc}
\toprule
\textbf{Signal} & \textbf{Frequency dep.} & \textbf{Time dep.} & \textbf{DFD?} \\
\midrule
DFD scalar field & None (universal) & Annual + secular & Yes \\
Dark photon & Mass-dependent & Stochastic & No \\
Axion & Mass-dependent & Coherent oscillation & No \\
Thermal drift & None & Irregular & No \\
\bottomrule
\end{tabular}
\end{center}

\subsubsection*{3. Experimental Approach}

\paragraph{3.1 Parasitic mode.} Use existing Dark SRF frequency monitoring data ($\geq 1$ year baseline). No modification to experimental setup. Analysis-only.

\paragraph{3.2 Search strategy.}
\begin{enumerate}
\item \textbf{Annual modulation:} Fourier analysis of cavity frequency time series at $f = 1/\text{year}$. Expected amplitude: $\delta\nu/\nu \sim 10^{-11}$. Signal-to-noise depends on frequency stability over 1 year.
\item \textbf{Secular drift:} Linear fit to long-baseline frequency data. Expected slope: $\dot\nu/\nu \sim 10^{-18}$ yr$^{-1}$. Requires $\geq 2$ years of data.
\item \textbf{Multi-cavity cross-check:} Simultaneous monitoring of two cavities at different frequencies. DFD signal is frequency-independent; systematic noise is not. Correlation analysis distinguishes signal from systematic.
\end{enumerate}

\paragraph{3.3 Systematics control.}
\begin{itemize}
\item Temperature: Already controlled at mK level in SQMS cryostats. Thermal expansion coefficient of Nb well-characterised.
\item Pressure: Cryostat provides vacuum isolation. No atmospheric coupling.
\item Magnetic field: Shielded to $< 1$ $\mu$T. Field-dependent frequency shift is $\sim 10^{-12}$ Hz/$\mu$T, negligible.
\item Helium pressure: He bath pressure variations at $\sim 10^{-3}$ mbar level produce $\delta\nu/\nu \sim 10^{-12}$. Can be correlated and subtracted.
\end{itemize}

\subsubsection*{4. Resources}

\begin{center}
\begin{tabular}{lrl}
\toprule
\textbf{Item} & \textbf{Cost} & \textbf{Notes} \\
\midrule
Postdoc (0.5 FTE, 12 months) & \$65,000 & Data analysis, paper writing \\
Computing (HPC) & \$5,000 & $\sim 10^4$ CPU-hours \\
Travel & \$5,000 & 1 conference + 2 SQMS visits \\
\midrule
\textbf{Total} & \textbf{\$75,000} & From SQMS 2.0 budget \\
\bottomrule
\end{tabular}
\end{center}

Upper estimate \$100k if full FTE is required.

\subsubsection*{5. Expected Outcomes}

\begin{enumerate}
\item \textbf{Detection:} Annual modulation at predicted amplitude and phase would constitute first laboratory evidence for DFD scalar coupling. Result: 1 PRL-class paper.
\item \textbf{Null result:} Upper bound on $\psi$-EM coupling at $< 10^{-10}$ level, improving on current constraints by $\sim 2$ orders of magnitude. Result: 1 PRD paper.
\item \textbf{Either outcome:} Demonstrates novel use of SRF infrastructure for fundamental physics beyond the axion/dark photon programme. Strengthens the SQMS science case.
\end{enumerate}

\subsubsection*{6. Timeline}

\begin{center}
\begin{tabular}{ll}
\toprule
\textbf{Date} & \textbf{Milestone} \\
\midrule
Q2 2026 & Submit white paper to SQMS Internal Review Committee \\
Q3 2026 & Present at SQMS 2.0 kickoff meeting \\
Q4 2026 & Begin data analysis (if approved) \\
Q2 2027 & Preliminary results \\
Q3 2027 & Paper submission \\
\bottomrule
\end{tabular}
\end{center}

\subsubsection*{7. Contact}

This proposal is submitted by G.T.~Alcock. We seek endorsement from a member of the SQMS Dark SRF team (suggested: A.~Romanenko, Dark SRF PI) as co-proposer.

\noindent\rule{\textwidth}{0.5pt}

\subsection{SQMS Submission Path --- Updated}

\begin{enumerate}
\item \textbf{April 2026:} Format as SQMS-NOTE. Assign document number.
\item \textbf{April--May 2026:} Circulate to A.~Romanenko (Dark SRF PI) and R.~Pilipenko (SQMS Deputy Director).
\item \textbf{Q3 2026:} Present at SQMS 2.0 kickoff meeting. The \$125M renewal provides a favourable environment for parasitic proposals.
\item \textbf{Q3--Q4 2026:} If endorsed, formal submission to Internal Review Committee.
\item \textbf{Q4 2026:} Begin analysis (0.5 FTE postdoc hire or existing SQMS postdoc allocation).
\end{enumerate}

\textbf{Agent 8 Status: \GREEN.} SQMS white paper in final form. Ready for formatting as SQMS-NOTE and circulation. \$125M SQMS 2.0 renewal creates favourable environment.

%=============================================================================
\section{Patent Agents --- Combined i31 Assessment}
%=============================================================================

\begin{center}
\begin{tabular}{clll}
\toprule
\textbf{Agent} & \textbf{Patents} & \textbf{Status} & \textbf{Action} \\
\midrule
6 & P1--P4 & \RED (dead) & Final disposition filed. Monitoring ceased \\
7 & P5--P7 & \GREEN & IR1 Q3 2026; ET site H2 2026; CE shortlist 2028 \\
8 & P8--P11 & \GREEN & SQMS white paper final. Submit Q2 2026 \\
\bottomrule
\end{tabular}
\end{center}

\subsection{TOP 5 FINDINGS --- Patent Agents i31}

\begin{enumerate}
\item \textbf{SQMS 2.0 funded at \$125M over 5 years.} The Dark SRF programme continues as flagship. Parasitic DFD search is feasible within this budget at \$75--100k.

\item \textbf{ET site decision expected H2 2026.} Double-L configuration (Sardinia + second site) improves DFD testability through better sky localisation and polarisation measurement.

\item \textbf{IR1 observing run starts September--October 2026.} Six-month run with both LIGO detectors. DFD $c_T = c$ prediction continues as null test.

\item \textbf{SQMS white paper finalised.} Complete 2-page document with executive summary, science case, experimental approach, resources, timeline, and expected outcomes. Ready for SQMS-NOTE formatting.

\item \textbf{P1--P4 closed permanently.} No experimental path exists for local DFD detection given the screening mechanism. Resources redirected to cosmological and astrophysical tests.
\end{enumerate}

\end{document}
